\ \\ 
\textbf{Утверждение} Элемент $\alpha \in L$ — алгебраический над $F$ $\Longleftrightarrow$ расширение $F(\alpha)$ конечно, и в
этом случае $F(\alpha) = F[\alpha] \sim= F[x]/(m \alpha)$. \\
\Proof \\
$\Rightarrow$ Зафиксируем $m_{\alpha}$ — минимальный многочлен элемента $\alpha$, и рассмотрим гомоморфизм
$\phi$ : $F$[$x$] $\rightarrow$ F[$\alpha$] такой, что $\forall f \in F$[$x$] : $f \rightarrow f(\alpha)$. Тогда $\phi$ сюръективен и $Ker \phi$ = ($m_{\alpha}$),
значит, по теореме о гомоморфизме, $F$[$\alpha$] $\sim =$ $F$[$x$]/($m_{\alpha}$). Следовательно, $F$[$\alpha$] — поле,
поэтому $F$[$\alpha$] = $Quot F$[$\alpha$] = $F$($\alpha$) и [$F$($\alpha$) : $F$] = deg $m_{\alpha}$.\\
$\Leftarrow$ В силу конечности расширения $F$($\alpha$), система (1, $\alpha$, $\alpha^2$, . . .) линейно зависима, поэтому
существует такой нетривиальный многочлен $f \in F[x] \backslash \{0\}, что f(\alpha) = 0.$
\EndProof \\ 
\textbf{Утверждение} \\ 
Пусть $L \supset F$ — конечное расширение. Тогда каждый элемент $\alpha \in L$
является алгебраическим над $F$.
\\
\Proof  Расширение $L$ конечно, причем $L \supset F(\alpha) \supset F$, поэтому расширение $F$($\alpha$) также конечно. \EndProof